\chapter{Experimenteller M12-Strömungskern}
\label{api:cfd}
Die Module unter \texttt{fem\_toolbox.cfd} sind Forschungs-APIs.
MAC bleibt die unabhängige laminare Referenz; Quad-FV ergänzt körperangepasste
ebene Netze. SA-noft2 ist als Forschungsentwicklung ergänzt; FSI ist noch offen.
\texttt{FluidMesh2D} prüft SI-Punkte, konvexe Quad4-Zellen und Line2-Randgruppen.
\texttt{FluidDomain.from\_file} liest ein Schema-1-JSON mit expliziter Breite,
positiver Dichte/Viskosität und genau einer physikalischen Randbedingung je Face.
Einlass, Wand, Druckauslass, Symmetrie und kongruente periodische Paare sind
getrennt; ohne Auslass ist genau eine Druckreferenzzelle nötig.

\texttt{solve\_quad\_flow} benötigt positive Geschwindigkeits- und Längenskalen.
Geschwindigkeit hat Form $(n_c,2)$ in m/s, Druck $(n_c)$ in Pa und Massenfluss
$(n_f)$ in kg/s; die Normalen zeigen aus der Eigentümerzelle.
Variable effektive Viskosität ist als Zell- und Randfeld zusammen zu liefern.
Die vollen Newton-Spannungen bleiben erhalten. Ergebnisse und Provenienz sind
unveränderlich. Rohresiduen, globale Masse/Impuls und diskrete Leistung werden
nach dem Update frisch rekonstruiert. Kleine Iterandenänderung genügt nicht.
Ungültige Geometrie, Rückströmung, Ressourcenlimit und Nichtkonvergenz führen
zu konkreten Fehlern; kein stiller Ersatz durch einen einfacheren Solver.

Die exakten Signaturen folgen im Inventar. Manuelle Detailverträge:
\nolinkurl{Doc/70_STATIONARY_FLOW_CORE.md},
\nolinkurl{Doc/72_FLUID_MESH_ADAPTER.md},
\nolinkurl{Doc/73_QUAD_FV_OPERATORS.md} und
\nolinkurl{Doc/74_QUAD_FV_SIMPLE.md}.

\section{SA-noft2 und gemeinsame RANS-Iteration}
\texttt{sa\_model.sa\_terms} wertet die vollständigen Punktterme mit
Allmaras-S-Modifikation aus. \texttt{sa\_transport.solve\_sa\_transport}
löst die skalare Gleichung bei vorgegebenem konservativem Massenstrom;
explizite MMS-Rand-/Wandabstandsfelder sind als Verifikationsdaten markiert.
\texttt{solve\_rans} koppelt U, p und a mit gemeinsamer Positivitäts-Schrittsuche.
Alle Quellen, variable Spannungen und Residuen werden am akzeptierten Zustand
neu rekonstruiert. Wandauflösung wird mit tatsächlicher Schubspannung geprüft.
Felder und Fehlerhistorien sind unveränderlich. SI-Formen, Fehlerzustände,
vorab gesetzte Gates und getrennte Code-/Modellabnahme stehen in
\nolinkurl{Doc/88_SA_NOFT2_IMPLEMENTATION.md}. Keine FSI- oder allgemeine
Turbulenzmodellfreigabe folgt allein aus einem konvergierten Lauf.
